\SourceLeaf{333}%
\Article{32}
Schließlich ergeben sich für die sieben Bestandteile von \(H\) die folgenden Werte
\[
\begin{aligned}
1)\quad&-\frac{kk\cos T^2}{aacc}\mathord{.}ss\\
2)\quad&-\frac{kk\cos T^2}{12aa(1-ee)cc}\mathord{.}(1-2eess+ees^4)\\
3)\quad&+\frac{kk\sin T^2}{12aacc}\mathord{.}ss\\
4)\quad&+\frac{eekk\cos T^2}{4aa(1-ee)}\mathord{.}ss\\
5)\quad&-\frac{kk\sin T^2}{24aacc}\mathord{.}ss\\
6)\quad&+\frac{kk\cos T^2}{12aacc}\mathord{.}ss\\
7)\quad&+\frac{kk\cos T^2}{24aa(1-ee)cc}\mathord{.}(1+ss-2ees^4)
\end{aligned}
\]
Die Glieder 1 und 6 heben sich gegenseitig auf. Fasst man die übrigen zusammen, so erhält man
\[
H=\frac{kk\sin T^2}{24aacc}\mathord{.}ss-\frac{kk\cos T^2}{24aa(1-ee)}\mathord{.}(1-10eess)
\]
woraus sich gemäß \(l=\lambda(1+Hrr)\) ergibt
\[
l=\lambda(1+\tfrac1{24}\tau\tau-\frac{1-ee}{24k^4}\mathord{.}(1-10eess)\beta\beta)\quad .\ .\ .\ .\ .\ .\quad(11)
\]

Die Formeln 9, 10, 11, die die Lösung unserer Aufgabe enthalten, unterscheiden sich von den Formeln III, II, I (Art. 25, 24, 23) nur darin, dass in den ersteren innerhalb der Klammern \(\tau\) und \(\beta\) dort stehen, wo in den letzteren \(t\) und \(b\) stehen. Wie man leicht sieht, entstehen dadurch in den Endergebnissen nur Unterschiede fünfter Ordnung. Da die ersteren Formeln, wie aus ihrer Herleitung hervorgeht, bis zur fünften Ordnung (diese ausgenommen) genau sind, ist bewiesen, dass auch die nach der ersten Methode gefundenen Formeln I, II, III (Art. 23---25) dieselbe Genauigkeit besitzen.

\Needspace{6\baselineskip}
\Article{33}
Für die numerische Berechnung wird man die Formeln 9, 10, 11 vorzugsweise in die folgende logarithmische Form bringen, bei der der Grad der Genauigkeit offensichtlich unverändert bleibt. Dabei bezeichnet \(M\) den Modulus des gewählten Logarithmensystems.
